\documentclass[11pt,a4paper]{article}
\usepackage[T1]{fontenc}
\usepackage[utf8]{inputenc}
\usepackage[margin=1in]{geometry}
\usepackage{lmodern}
\usepackage{amsmath,amssymb,graphicx,booktabs,tabularx,array}
\usepackage{microtype}
\usepackage{caption}
\usepackage{placeins}
\usepackage[hidelinks]{hyperref}
\usepackage{float}

\captionsetup{font=small,labelfont=bf}
\setlength{\emergencystretch}{2em}
\newcolumntype{Y}{>{\raggedright\arraybackslash}X}

\newcommand{\figpath}{../figures/Whakaari}

\newcommand{\pending}[1]{\par\smallskip\noindent\textit{To complete: #1}\par\smallskip}

\title{Ordinal Pattern Analysis of RSAM Time Series at Whakaari, 2014--2019}
\author{Rasika Dilhani}
\date{\today}

\begin{document}
	\maketitle
	\begin{abstract}
		This study investigates ordinal pattern features for retrospective detection of unusual seismic-amplitude behavior at Whakaari, New Zealand. RSAM records from WIZ and WSRZ cover the intended period July 2014-December 2019. Seven descriptors from thirteen features are used for DBSCAN anomaly detection and weighted radial-basis SVM classification against three Volcanic Alert Level classes. The corrected modeling datasets contain 335 and 333 weekly windows, respectively, including three eruption-labeled windows at each station. DBSCAN flags 42 WIZ and 46 WSRZ windows, all outside the eruption class. Repeated three-fold SVM validation gives mean one-versus-rest balanced accuracies of 0.538 and 0.642. A separate out-of-fold pass detects zero of three eruption windows at WIZ and one of three at WSRZ. These exploratory results identify unusual feature-space observations, but do not establish reliable eruption recognition or forecasting.
	\end{abstract}
	
\section{Introduction}
Real-time Seismic Amplitude Measurement (RSAM) summarizes seismic amplitude over specified intervals and supports volcano monitoring~\cite{Endo1991}. Amplitude summaries provide a useful overview, while changes in the temporal ordering of observations may supply complementary information. Ordinal-pattern analysis represents this ordering through distributions of ranked observation sequences~\cite{Bandt2002}.

This study examines whether ordinal features identify unusual weekly RSAM windows at Whakaari and how these windows relate to Volcanic Alert Level (VAL) records. We analyze WIZ and WSRZ separately and compare them over a common study period. DBSCAN provides an unsupervised anomaly flag; a weighted support vector machine (SVM) addresses supervised activity classification. The objective is retrospective characterization and detection, with forecasting claims requiring a separate evaluation using only information available before each event.
	
\section{Study Area and Data Overview}\label{sec:data}
\subsection{Records and Scope}
The study concerns Whakaari, New Zealand, and RSAM records identified by stations WIZ and WSRZ. The intended analysis interval is 1 July 2014-31 December 2019. Figure axes are labeled in New Zealand local time. 	
	
\subsection{Long-Term RSAM Records}
Figures~\ref{fig:full_wiz} and~\ref{fig:full_wsrz} introduce the full records before the feature analysis. Both show a relatively low baseline punctuated by transient peaks, with differences in amplitude and timing between stations. 
	
\begin{figure}[H]
	\centering
	\includegraphics[width=\linewidth]{\figpath/WIZ_2014to2019_full_rsam_highlighted.pdf}
	\caption{WIZ RSAM record for the 2014-2019 interval.}
	\label{fig:full_wiz}
\end{figure}	
	
\begin{figure}[H]
	\centering
	\includegraphics[width=\linewidth]{\figpath/WSRZ_2014to2019_full_rsam_highlighted.pdf}
	\caption{WSRZ RSAM record for the 2014-2019 interval.}
	\label{fig:full_wsrz}
\end{figure}	
	
	
\subsection{Selected Short-Interval Views}
Figures~\ref{fig:apr2016}-\ref{fig:dec2019interval} compare the two stations around the dates selected in the supplied plots. Green represents WIZ and purple represents WSRZ. 
	
\begin{figure}[H]
	\centering
	\includegraphics[width=\linewidth]{\figpath/zoom_28_Apr_2016_eruption.pdf}
	\caption{RSAM comparison around the 28 April 2016 annotation.}
	\label{fig:apr2016}
\end{figure}

\begin{figure}[H]
	\centering
	\includegraphics[width=\linewidth]{\figpath/zoom_13_15_Sep_2016_eruption.pdf}
	\caption{RSAM comparison around the 13-15 September 2016 annotation. The isolated WIZ peak precedes the highlighted interval.}
	\label{fig:sep2016}
\end{figure}

\FloatBarrier
\begin{figure}[H]
	\centering
	\includegraphics[width=\linewidth]{\figpath/zoom_9_Dec_2019_explosive_eruption.pdf}
	\caption{RSAM comparison around 9 December 2019.}
	\label{fig:dec2019onset}
\end{figure}

\begin{figure}[H]
	\centering
	\includegraphics[width=\linewidth]{\figpath/zoom_9_12_Dec_2019_continued_eruptive_activity.pdf}
	\caption{RSAM comparison with the 9-12 December interval.}
	\label{fig:dec2019interval}
\end{figure}	
	
	
\section{Literature Review}
	
\subsection{Information-Theoretic Measures in Volcanic Seismology}
	
Information-theoretic measures quantitatively describe the variability and statistical structure of seismic signals. In volcanic monitoring, researchers have investigated these measures as complementary tools for identifying changes in seismic activity associated with eruptive processes. Telesca et al.~\cite{Telesca2014} applied the Fisher--Shannon method to continuous seismic recordings collected before and during the 2011-2012 El Hierro eruption in the Canary Islands. Their work illustrates how complementary information measures can characterize the evolution of volcanic seismic signals.
	
Rey-Devesa et al.~\cite{ReyDevesa2023} investigated Shannon entropy in continuous seismic recordings from Volc\'an de Colima, Mexico. They reported decreases in entropy preceding two major eruptive episodes and examined entropy minima as indicators of smaller explosions, supported by visual monitoring observations. These findings suggest that entropy can provide information relevant to tracking volcanic activity. However, entropy variations alone do not establish a unique physical mechanism or demonstrate deterministic conduit behavior. Furthermore, Shannon entropy derived from seismic waveform statistics should be distinguished from permutation entropy, which describes the probability distribution of temporal ordinal patterns.
	
\subsection{Ordinal Patterns and Permutation Entropy}
	
Bandt and Pompe~\cite{Bandt2002} introduced permutation entropy as a method for characterizing time series through the relative ordering of observations. For a selected embedding dimension $D$ and time delay $\tau$, sequences of observations are represented by ordinal patterns. The relative frequencies of these patterns define a probability distribution, from which permutation entropy is calculated. This representation captures aspects of temporal organization without requiring a specified parametric distribution for the signal amplitudes. Although the method is computationally efficient and can be useful for noisy observations, its performance depends on the embedding parameters, the number of available observations, and how tied values are treated. Its application to non-stationary signals also requires analysis windows appropriate to the timescales under investigation.
	
Statistical complexity complements entropy by capturing a distribution's departure from a reference state. Lamberti et al.~\cite{lamberti2004intensive} developed an intensive statistical complexity measure using Jensen--Shannon divergence to quantify disequilibrium. In this framework, normalized Shannon entropy $H$ and normalized disequilibrium $Q_J$ are combined as $C=H Q_J$. When calculated from ordinal-pattern distributions, these measures provide a joint representation of temporal pattern diversity and departure from a uniform distribution in the $H \times C$ plane.
	
\subsection{State Transitions in Geysers and Implications for Volcanic Monitoring}
	
Sudibyo et al.~\cite{Sudibyo2022} investigated permutation entropy in seismic recordings from Strokkur geyser, Iceland. Their analysis of 63 eruptive cycles identified recurring entropy variations associated with different eruption-cycle phases. They reported a systematic increase in permutation entropy during the final approximately 15 s before eruption. They found that permutation entropy outperformed seismic signal strength or quiescence for forecasting eruptions in the studied setting. Their work also examined how parameter selection and station distance affected the resulting entropy measurements.
	
These findings provide empirical motivation for investigating ordinal pattern measures as indicators of changes in volcanic and hydrothermal activity. However, transferring results from Strokkur to Whakaari requires separate evaluation. Repeated geyser eruptions observed through high-frequency seismic recordings differ from volcanic activity examined through aggregated RSAM time series over weekly windows. Consequently, the direction, timing, and physical interpretation of entropy changes cannot be assumed to remain consistent across these systems.
	
In the present study, the Strokkur investigation is relevant because it shows that temporal ordinal structure can contain information about changes in an eruptive system. This supports examining ordinal feature trajectories in Whakaari records while distinguishing retrospective anomaly detection from eruption forecasting. A window containing eruptive or post-eruptive observations may help characterize an eruption-related change, but it cannot establish a pre-eruptive signal.
	
\subsection{Motivation for Multi-Station Feature Analysis}
The reviewed studies motivate examining information measures in long-term RSAM records, but do not directly establish the performance of the present combination of ordinal features at Whakaari. Entropy, complexity, disequilibrium, and variance descriptors may overlap. Feature selection is therefore used to seek a compact representation for clustering and classification. Comparing WIZ and WSRZ can reveal whether detected changes are shared across stations, while retaining station-specific differences in data quality and signal characteristics.
	
\section{Methodology} 

\subsection{Preprocessing and Weekly Windows} 

The feature-generation description specifies 10-minute RSAM observations and seven-day windows. The supplied clustering script reads precomputed weekly features and does not construct windows; therefore, this script does not verify the step size. A complete window on a regular UTC grid contains $N=7\times144=1008$ observations. The window duration and step are separate settings: a seven-day duration does not imply a seven-day step. 

For the full intended interval, a daily step would produce about 2000 candidate windows before exclusions, whereas the corrected modeling output contains 335 WIZ and 333 WSRZ windows. The difference depends on the feature-generation settings and exclusion counts, not an assumed daily step. Seven local calendar days can also differ from 168 hours at daylight-saving transitions; the time convention must be explicit. 


\subsection{Ordinal Pattern Distribution}

For a regular time series $X=\{x_t\}_{t=1}^{N}$, ordinal patterns are constructed from the embedding vectors
\begin{equation}
	\mathbf{x}_t =
	(x_t,x_{t+\tau},\ldots,x_{t+(D-1)\tau}),
\end{equation}
where $D$ denotes the embedding dimension and $\tau$ is the embedding delay. Each vector is mapped to the permutation that describes the relative ordering of its components.

In this study, the embedding parameters are $D=5$ and $\tau=1$ observation. Therefore, the number of possible ordinal patterns is
\begin{equation}
	k=D!=5!=120.
\end{equation}
For a complete seven-day window containing $N=1008$ observations, the number of available embedding vectors is
\begin{equation}
	M=N-(D-1)\tau=1004.
\end{equation}

Let $n_\ell$ denote the number of occurrences of ordinal pattern $\ell$. The empirical ordinal-pattern probability distribution is then
\begin{equation}
	\widehat{p}_\ell=\frac{n_\ell}{M},
	\qquad
	\ell=1,\ldots,k,
\end{equation}
and
\begin{equation}
	\widehat{\mathbf{P}}
	=
	\{\widehat{p}_\ell\}_{\ell=1}^{k},
	\qquad
	\sum_{\ell=1}^{k}\widehat{p}_\ell=1.
\end{equation}

If embedding vectors containing invalid or missing observations are excluded, $M$ denotes the number of valid vectors actually used to estimate the distribution. A consistent tie-breaking rule is applied when equal values occur within an embedding vector. Because the discrete Fisher Information Measure depends on differences between adjacent pattern probabilities, a fixed ordering of the ordinal patterns is also retained throughout the analysis.

\subsection{Information-Theoretic Quantifiers}

For each empirical ordinal-pattern probability distribution
$\widehat{\mathbf{P}}=\{\widehat{p}_\ell\}_{\ell=1}^{k}$,
a collection of global and local information-theoretic measures is calculated. These measures characterize different properties of the temporal dynamics represented by the ordinal-pattern distribution. Entropy-based measures describe global uncertainty and disorder, whereas the Fisher Information Measure describes local variation among neighbouring probabilities.

\subsubsection{Permutation Entropy}

Shannon entropy provides a measure of the uncertainty contained in the ordinal-pattern probability distribution. The Shannon entropy of
$\widehat{\mathbf{P}}$ is
\begin{equation}
	S(\widehat{\mathbf{P}})
	=
	-\sum_{\ell=1}^{k}
	\widehat{p}_\ell\log \widehat{p}_\ell,
\end{equation}
where terms with $\widehat{p}_\ell=0$ are defined to contribute zero.

The maximum Shannon entropy occurs for the uniform distribution and is equal to $\log k$. The normalized Shannon entropy, commonly referred to as permutation entropy in the ordinal-pattern framework, is therefore defined as
\begin{equation}
	H_{S}(\widehat{\mathbf{P}})
	=
	-\frac{1}{\log k}
	\sum_{\ell=1}^{k}
	\widehat{p}_\ell\log \widehat{p}_\ell.
\end{equation}
Consequently,
\begin{equation}
	0\leq H_{S}(\widehat{\mathbf{P}})\leq1.
\end{equation}

Low permutation entropy indicates that only a small number of ordinal patterns dominate the distribution, corresponding to a comparatively regular or ordered signal. Values close to one indicate a distribution approaching uniformity and therefore greater disorder or unpredictability. Entropy alone, however, cannot distinguish complete randomness from more structured forms of dynamical organization. For this reason, additional information and complexity measures are considered.


\subsubsection{Tsallis Entropy}

For $q\in\mathbb{R}\setminus\{1\}$, the Tsallis entropy is defined as
\begin{equation}
	S_{T}^{(q)}(\widehat{\mathbf{P}})
	=
	\frac{1}{q-1}
	\left(
	1-\sum_{\ell=1}^{k}
	\widehat{p}_\ell^{\,q}
	\right).
\end{equation}
The Shannon entropy is recovered in the limit $q\rightarrow1$.

The maximum Tsallis entropy occurs for the uniform distribution and is
\begin{equation}
	S_{T,\max}^{(q)}
	=
	\frac{1-k^{1-q}}{q-1}.
\end{equation}
Hence, the normalized Tsallis entropy used in the analysis is
\begin{equation}
	H_{T}^{(q)}(\widehat{\mathbf{P}})
	=
	\frac{
		S_{T}^{(q)}(\widehat{\mathbf{P}})
	}{
		S_{T,\max}^{(q)}
	}.
\end{equation}

The parameter $q$ controls the relative contribution of common and rare ordinal patterns and therefore provides an alternative characterization of the ordinal-pattern distribution to Shannon entropy.

\subsubsection{R\'enyi Entropy}

For $q>0$, $q\neq1$, the R\'enyi entropy is defined as
\begin{equation}
	S_{R}^{(q)}(\widehat{\mathbf{P}})
	=
	\frac{1}{1-q}
	\log
	\left(
	\sum_{\ell=1}^{k}
	\widehat{p}_\ell^{\,q}
	\right).
\end{equation}
As with Tsallis entropy, Shannon entropy is obtained in the limit $q\rightarrow1$.

The maximum R\'enyi entropy is $\log k$, obtained for the uniform distribution. The normalized R\'enyi entropy is therefore
\begin{equation}
	H_{R}^{(q)}(\widehat{\mathbf{P}})
	=
	\frac{
		S_{R}^{(q)}(\widehat{\mathbf{P}})
	}{
		\log k
	}.
\end{equation}

Different values of $q$ change the relative emphasis placed on ordinal patterns with high and low probabilities and provide complementary descriptions of the underlying temporal dynamics.


\subsubsection{Fisher Information Measure}

Entropy measures characterize global properties of the ordinal-pattern distribution. To complement these measures with a descriptor that is sensitive to local changes in the probability distribution, we adopt the discrete Fisher Information Measure proposed by Rey et al.~\cite{Rey2024Asymptotic}:
\begin{equation}
	F(\widehat{\mathbf{P}})
	=
	4
	\sum_{\ell=1}^{k-1}
	\left(
	\sqrt{\widehat{p}_{\ell+1}}
	-
	\sqrt{\widehat{p}_{\ell}}
	\right)^2.
\end{equation}

Unlike entropy, which depends on the overall distribution of probabilities, the Fisher Information Measure depends on differences between neighbouring probabilities under the adopted ordinal-pattern ordering. It therefore captures local gradients and fine-scale changes in the probability distribution that may not be evident from entropy measures alone.


\subsubsection{Disequilibrium}
\label{disequilibrium}

The entropy and Fisher Information measures characterize complementary properties of the ordinal-pattern distribution. However, an additional quantity is required to measure the departure of the observed distribution from complete randomness. In the ordinal-pattern framework, the natural equilibrium reference is the uniform distribution
\begin{equation}
	\mathbf{U}
	=
	\left(
	\frac{1}{k},
	\frac{1}{k},
	\ldots,
	\frac{1}{k}
	\right).
\end{equation}

Following the disequilibrium concept introduced by Lopez-Ruiz et al.~(\cite{lopez1995statistical}), the departure from the uniform distribution is quantified using the Jensen--Shannon divergence. Let
\begin{equation}
	S(\mathbf{P})
	=
	-\sum_{\ell=1}^{k}p_\ell\log p_\ell
\end{equation}
denote the unnormalized Shannon entropy. The Jensen--Shannon divergence between the empirical distribution $\widehat{\mathbf{P}}$ and the uniform distribution $\mathbf{U}$ is
\begin{equation}
	J(\widehat{\mathbf{P}},\mathbf{U})
	=
	S\left(
	\frac{\widehat{\mathbf{P}}+\mathbf{U}}{2}
	\right)
	-
	\frac{1}{2}S(\widehat{\mathbf{P}})
	-
	\frac{1}{2}S(\mathbf{U}),
\end{equation}
where
\begin{equation}
	S(\mathbf{U})=\log k.
\end{equation}

The maximum possible Jensen--Shannon divergence from the uniform distribution occurs when all probability mass is concentrated on a single ordinal pattern and is given by
\begin{equation}
	J_{\max}
	=
	\log 2
	+
	\frac{1}{2}\log k
	-
	\frac{k+1}{2k}\log(k+1).
\end{equation}

The normalized disequilibrium is therefore defined as
\begin{equation}
	Q(\widehat{\mathbf{P}})
	=
	\frac{
		J(\widehat{\mathbf{P}},\mathbf{U})
	}{
		J_{\max}
	}.
\end{equation}
Hence,
\begin{equation}
	0\leq Q(\widehat{\mathbf{P}})\leq1.
\end{equation}

For the uniform distribution,
\begin{equation}
	Q(\widehat{\mathbf{P}})=0,
\end{equation}
indicating complete equilibrium or maximal disorder. Conversely, $Q(\widehat{\mathbf{P}})=1$ when the entire probability mass is concentrated on a single ordinal pattern, corresponding to the maximum departure from the uniform distribution.

The same Jensen--Shannon-based disequilibrium is used throughout this study to construct the Shannon, Tsallis, R\'enyi, and Fisher-based statistical complexity measures.


\subsubsection{Statistical Complexity Measures}
\label{complexity-measures}

Entropy measures quantify global uncertainty, whereas Fisher information quantifies local changes in the ordinal-pattern distribution. Neither type of measure alone explicitly incorporates the departure of the observed distribution from the uniform distribution. Statistical complexity measures combine these complementary characteristics.

Following the framework of Lamberti et al.~\cite{lamberti2004intensive}, the statistical complexity is constructed by combining an information measure with the normalized disequilibrium $Q(\widehat{\mathbf{P}})$. This construction produces low complexity for distributions approaching complete disorder and also distinguishes highly concentrated distributions from distributions exhibiting intermediate levels of organization.

The Shannon statistical complexity is defined as
\begin{equation}
	C_{S}(\widehat{\mathbf{P}})
	=
	H_{S}(\widehat{\mathbf{P}})
	Q(\widehat{\mathbf{P}}).
\end{equation}

Similarly, the Tsallis statistical complexity is
\begin{equation}
	C_{T}^{(q)}(\widehat{\mathbf{P}})
	=
	H_{T}^{(q)}(\widehat{\mathbf{P}})
	Q(\widehat{\mathbf{P}}),
\end{equation}
where $H_{T}^{(q)}(\widehat{\mathbf{P}})$ is the normalized Tsallis entropy of order $q$.

The R\'enyi statistical complexity is defined as
\begin{equation}
	C_{R}^{(q)}(\widehat{\mathbf{P}})
	=
	H_{R}^{(q)}(\widehat{\mathbf{P}})
	Q(\widehat{\mathbf{P}}),
\end{equation}
where $H_{R}^{(q)}(\widehat{\mathbf{P}})$ denotes the normalized R\'enyi entropy of order $q$.

Finally, the Fisher-based complexity is defined by combining the Fisher Information Measure with the same normalized disequilibrium:
\begin{equation}
	C_{F}(\widehat{\mathbf{P}})
	=
	F(\widehat{\mathbf{P}})
	Q(\widehat{\mathbf{P}}).
\end{equation}

Together, the Shannon, Tsallis, and R\'enyi entropy-complexity measures characterize global properties of the ordinal-pattern distribution, whereas the Fisher Information Measure and Fisher-based complexity provide complementary sensitivity to local structure. These quantities are subsequently used as features to characterize temporal changes in the RSAM time series.

\subsection{Thirteen Features from ordinal patterns}
Table~\ref{tab:features} lists the thirteen features specified for this analysis. Here $H_{\mathrm{Fisher}}$ retains the analysis-file notation but denotes Fisher information, not an entropy.
	
\begin{table}[H]
	\centering
	\caption{Candidate features, grouped by type.}\label{tab:features}
	\begin{tabularx}{\linewidth}{lYr}
		\toprule
		Group & Features & Count\\
		\midrule
		Entropy/information & $H_{\mathrm{Shannon}}$, $H_{\mathrm{R\acute{e}nyi}}$, $H_{\mathrm{Tsallis}}$, $H_{\mathrm{Fisher}}$ & 4\\
		Complexity & $C_{\mathrm{Shannon}}$, $C_{\mathrm{R\acute{e}nyi}}$, $C_{\mathrm{Tsallis}}$, $C_{\mathrm{Fisher}}$ & 4\\
		Disequilibrium & Disequilibrium & 1\\
		Variance & $\mathrm{Var}(H_{\mathrm{Shannon}})$, $\mathrm{Var}(H_{\mathrm{R\acute{e}nyi}})$, $\mathrm{Var}(H_{\mathrm{Tsallis}})$, $\mathrm{Var}(C_{\mathrm{Shannon}})$ & 4\\
		\bottomrule
	\end{tabularx}
\end{table}

The four complexity measures combine their corresponding normalized information measure with the Jensen--Shannon-based normalized disequilibrium defined in Section~\ref{disequilibrium}.
	
\subsection{Seven-Feature Representation}
The seven features retains Fisher information, Fisher-based complexity, disequilibrium, and the four variance measures in Table~\ref{tab:features}. This defines a seven-dimensional vector for each window. 
	
\subsection{VAL Reference Classes}\label{sec:val}
	
Volcanic Alert Level (VAL) records were obtained from the GeoNet data repository~\cite{GeoNetVAL}. These records describe assessed volcanic status and are used as an external activity reference, rather than as direct annotations of individual seismic events. Each record includes the period start and end, the numeric VAL, the VAL system version, and the official activity description. Period boundaries are provided in both UTC and New Zealand local time (NZST or NZDT).
	
The analysis uses VAL System Version 3 records overlapping the study period. Table~\ref{tab:val_v3} presents the official six-level classification.

\begin{table}[H]
	\centering
	\caption{Official GeoNet Volcanic Alert Level System
		Version 3 definitions~\cite{GeoNetVAL}.}
	\label{tab:val_v3}
	\begin{tabularx}{\linewidth}{lY}
		\toprule
		VAL & Volcanic activity \\
		\midrule
		0 & No volcanic unrest \\
		1 & Minor volcanic unrest \\
		2 & Moderate to heightened volcanic unrest \\
		3 & Minor volcanic eruption \\
		4 & Moderate volcanic eruption \\
		5 & Major volcanic eruption \\
		\bottomrule
	\end{tabularx}
\end{table}	
	
For clustering comparisons and supervised classification, the relevant Version 3 activity descriptions are grouped into three classes (Table~\ref{tab:val_mapping}). Class 1 	corresponds to minor volcanic unrest and is termed ``Background unrest'' in the analysis code. This label does not imply an absence of volcanic unrest. Class 2 represents moderate to heightened volcanic unrest, while Class 3 combines minor and moderate volcanic eruption. This mapping does not assign VAL 0 and VAL 5 a class.
	
\begin{table}[H]
	\centering
	\caption{Three-class grouping used in the analysis.
		The class numbers are study-specific labels, distinct
		from the official VAL scale.}
	\label{tab:val_mapping}
	\begin{tabularx}{\linewidth}{llY}
		\toprule
		Class & VAL & Description \\
		\midrule
		1 & 1 & Minor volcanic unrest
		(background unrest in the analysis) \\
		2 & 2 & Moderate to heightened volcanic unrest \\
		3 & 3--4 & Minor or moderate volcanic eruption \\
		\bottomrule
	\end{tabularx}
\end{table}	
	
For each weekly window identifies overlapping VAL periods and assigns the highest mapped class. An overlap occurs when the VAL period starts strictly before the window ends and ends strictly after the window starts. Windows with non-finite values in the selected features or without a mapped activity label are excluded. A window assigned to Class 3 therefore overlaps an eruption-labeled VAL period.
	
Figures~\ref{fig:val_system_sankey} and~\ref{fig:val_activity_sankey} illustrate the relationships between VAL levels, their official descriptions, and the three analysis classes. Flow widths represent the number of recorded VAL periods in the filtered dataset, rather than their durations or the number of weekly windows.
	
\begin{figure}[H]
	\centering
	\includegraphics[
	width=\linewidth,
	height=0.71\textheight,	keepaspectratio]
	{\figpath/VAL_System_Description_Sankey.pdf}
	\caption{Relationships between VAL levels and official
		activity descriptions in the filtered Version 3 record.
		Flow widths represent the number of recorded VAL periods.}
	\label{fig:val_system_sankey}
\end{figure}
	
\begin{figure}[H]
	\centering
	\includegraphics[
	width=\linewidth,
	height=0.71\textheight,
	keepaspectratio]{\figpath/VAL_Activity_Class_Sankey.pdf}
	\caption{Mapping from VAL levels to the three analysis
	classes. Flow widths represent numbers of recorded VAL
	periods, rather than their durations.}
	\label{fig:val_activity_sankey}
\end{figure}
	
	
\subsection{DBSCAN Anomaly Detection}
DBSCAN groups nearby feature vectors and labels points that do not satisfy its density conditions as noise. We treat noise points as candidate anomalous windows. They are not automatically eruption precursors, and membership of a dense cluster does not establish a physically normal state.
	
DBSCAN is fitted to all available standardized feature vectors separately for each station, with Euclidean distance, $\varepsilon=1.1$, and $\mathrm{minPts}=14$. The script generates a nearest-neighbor distance plot but fixes $\varepsilon$ in advance; it does not automatically select parameters. We compare noise descriptively with all three VAL classes and, for binary metrics, Class 3 versus Classes 1-2.
	
\subsubsection{Radial Basis Function Support Vector Machine}

A radial basis function (RBF) support vector machine (SVM) with class weighting was used for the supervised classification analysis. To account for class imbalance, the weight assigned to class \(k\) was calculated separately for each station as
\begin{equation}
	w_k = \frac{N}{3n_k},
\end{equation}
where \(N\) is the total number of windows and \(n_k\) is the number of windows belonging to class \(k\). The three classes therefore contribute more evenly to the SVM fitting procedure despite differences in their sample sizes.

The original feature values were supplied directly to the SVM routine without prior manual standardization. Scaling was instead performed internally using the default scaling procedure of the SVM implementation. This differs from the DBSCAN and PCA analyses, for which the features were explicitly standardized using the full dataset before analysis.

Hyperparameter selection for the RBF SVM was performed using a grid search over the cost parameter \(C\) and the RBF kernel parameter \(\gamma\). The parameter grid contained 40 combinations:
\begin{align*}
	C &\in \{0.1,\,1,\,10,\,100,\,1000\},\\
	\gamma &\in \{0.001,\,0.005,\,0.01,\,0.05,\,0.1,\,0.5,\,1,\,5\}.
\end{align*}

Because the smallest class contained only three windows, the initially requested five-fold cross-validation was reduced to three-fold stratified cross-validation. Each of the 40 hyperparameter combinations was evaluated using ten repeated stratified cross-validation partitions. The same random seeds and corresponding partitions were used across all parameter combinations to ensure a consistent comparison between candidate models.

For each repeat, the performance metrics were first averaged across the three folds and were subsequently averaged across the ten repeats. Hyperparameter selection was based on the mean one-versus-rest balanced accuracy across the three classes:
\begin{equation}
	\overline{\mathrm{BA}}
	=
	\frac{1}{3}
	\sum_{k=1}^{3}
	\frac{
		\mathrm{Sensitivity}_k
		+
		\mathrm{Specificity}_k
	}{2}.
\end{equation}

The hyperparameter combination producing the largest value of \(\overline{\mathrm{BA}}\) was selected as the final SVM configuration. This definition differs from macro recall because both class-specific sensitivity and specificity contribute to the balanced accuracy. Macro F1 was also calculated, with undefined class-specific F1 values replaced by zero before averaging. The reported standard deviation of balanced accuracy represents the variation among the ten repeat-level mean balanced accuracies.

Following hyperparameter selection, an additional seeded three-fold cross-validation was performed using the selected parameters to obtain one out-of-fold prediction for every window. These predictions were used to construct the confusion matrices and SVM classification figures. Consequently, the performance calculated from these out-of-fold predictions is not expected to be identical to the mean performance obtained from the repeated cross-validation procedure used for hyperparameter selection.

Since hyperparameter selection and performance assessment were conducted using the same available dataset, the reported validation results should not be interpreted as performance on an independent external test set. In addition, because some source windows overlap in time, observations assigned to different random folds may not be fully independent. A blocked or event-separated validation strategy would provide a more stringent assessment of model generalization under temporal dependence and is therefore identified as an important direction for future validation.
	
\section{Results}\label{sec:results}
\subsection{Data Coverage and VAL Class Counts}
The corrected output contains 335 WIZ and 333 WSRZ windows. Their start-date range is 1 July 2014-26 December 2019 in the code's local wall-clock representation. Table~\ref{tab:label_counts} shows the corrected counts. The code calls Class 1 ``Background unrest''; this corresponds to the VAL description ``Minor volcanic unrest'', not the absence of unrest.
	
\begin{table}[H]
	\centering\small
	\caption{Corrected weekly class counts.}\label{tab:label_counts}
	\begin{tabularx}{\linewidth}{Yrrrr}
		\toprule
		Station & Class 1 & Class 2 & Class 3 & Total\\
		\midrule
		WIZ & 320 & 12 & 3 & 335\\
		WSRZ & 318 & 12 & 3 & 333\\
		\bottomrule
	\end{tabularx}
\end{table}
	
\subsection{Entropy-Complexity Planes}
Figures~\ref{fig:hc_wiz} and~\ref{fig:hc_wsrz} retain the four supplied panels for each station. Shannon, Tsallis, and R\'enyi panels show concentrated curved relationships, while the Fisher panels show a different increasing association. Several labeled windows lie towards the edges of these distributions. WSRZ weeks 146--151 recur among the labeled points, and week 151 is labeled in several WIZ panels. Match week indices to dates before attributing event associations; they do not identify eruption windows. Colors encode week index rather than VAL class. The Fisher panel's horizontal coordinate denotes Fisher information despite the shared $H$ axis notation.

\begin{figure}[H]
	\centering
	\includegraphics[width=\linewidth,height=0.71\textheight,keepaspectratio]{\figpath/HC_Plane_WIZ_D5.pdf}
	\caption{WIZ ordinal-feature planes for $D=5$, with Shannon, Tsallis, R'enyi, and Fisher panels. Colors and selected week labels are retained from the supplied figure.}\label{fig:hc_wiz}
\end{figure}

\begin{figure}[H]
	\centering
	\includegraphics[width=\linewidth,height=0.71\textheight,keepaspectratio]{\figpath/HC_Plane_WSRZ_D5.pdf}
	\caption{WSRZ ordinal-feature planes for $D=5$. The panels retain their original scales, weak colors, labels, and guide lines.}\label{fig:hc_wsrz}
\end{figure}
	
Table~\ref{tab:hc_extremes_both_stations} compares the selected extreme windows for WIZ and WSRZ. Week 151, covering 17--24 December 2016, is selected in the Shannon, R\'enyi, and Tsallis planes at both stations. Among the
listed windows, it has the lowest entropy and highest complexity for these three measures at each station. The selected Fisher windows differ between stations: weeks 15, 68, 255, and 330 at WIZ, and weeks 116, 126, 147, and 151 at WSRZ. These selections describe extremes in the feature planes. 

\begin{table}[H]
	\centering
	\small
	\setlength{\tabcolsep}{4pt}
	\caption{Selected extreme windows for WIZ and WSRZ in the
		ordinal-feature planes ($D=5$).}
	\label{tab:hc_extremes_both_stations}
	\begin{tabular}{llrllrr}
		\toprule
		Measure & Station & Week & Start date & End date & $H$ & $C$ \\
		\midrule
		Shannon & WIZ & 91 & 2015-12-23 & 2015-12-30 & 0.91698 & 0.12138 \\
		Shannon & WIZ & 99 & 2016-02-09 & 2016-02-16 & 0.91937 & 0.12436 \\
		Shannon & WIZ & 151 & 2016-12-17 & 2016-12-24 & 0.89072 & 0.13647 \\
		Shannon & WIZ & 255 & 2018-09-02 & 2018-09-09 & 0.92004 & 0.12722 \\
		\addlinespace[2pt]
		Shannon & WSRZ & 146 & 2016-11-17 & 2016-11-24 & 0.88324 & 0.16222 \\
		Shannon & WSRZ & 147 & 2016-11-23 & 2016-11-30 & 0.87397 & 0.17108 \\
		Shannon & WSRZ & 150 & 2016-12-11 & 2016-12-18 & 0.88377 & 0.15582 \\
		Shannon & WSRZ & 151 & 2016-12-17 & 2016-12-24 & 0.85154 & 0.18292 \\
		\midrule
		Tsallis & WIZ & 91 & 2015-12-23 & 2015-12-30 & 0.96421 & 0.12319 \\
		Tsallis & WIZ & 99 & 2016-02-09 & 2016-02-16 & 0.96734 & 0.12629 \\
		Tsallis & WIZ & 104 & 2016-03-10 & 2016-03-17 & 0.96637 & 0.12472 \\
		Tsallis & WIZ & 151 & 2016-12-17 & 2016-12-24 & 0.94477 & 0.13971 \\
		\addlinespace[2pt]
		Tsallis & WSRZ & 146 & 2016-11-17 & 2016-11-24 & 0.94808 & 0.16807 \\
		Tsallis & WSRZ & 147 & 2016-11-23 & 2016-11-30 & 0.94233 & 0.17804 \\
		Tsallis & WSRZ & 150 & 2016-12-11 & 2016-12-18 & 0.94699 & 0.16115 \\
		Tsallis & WSRZ & 151 & 2016-12-17 & 2016-12-24 & 0.92604 & 0.19200 \\
		\midrule
		R\'enyi & WIZ & 91 & 2015-12-23 & 2015-12-30 & 0.87270 & 0.11150 \\
		R\'enyi & WIZ & 104 & 2016-03-10 & 2016-03-17 & 0.87938 & 0.11349 \\
		R\'enyi & WIZ & 151 & 2016-12-17 & 2016-12-24 & 0.81697 & 0.12081 \\
		R\'enyi & WIZ & 255 & 2018-09-02 & 2018-09-09 & 0.88691 & 0.11837 \\
		\addlinespace[2pt]
		R\'enyi & WSRZ & 146 & 2016-11-17 & 2016-11-24 & 0.82596 & 0.14642 \\
		R\'enyi & WSRZ & 147 & 2016-11-23 & 2016-11-30 & 0.81047 & 0.15313 \\
		R\'enyi & WSRZ & 150 & 2016-12-11 & 2016-12-18 & 0.82296 & 0.14005 \\
		R\'enyi & WSRZ & 151 & 2016-12-17 & 2016-12-24 & 0.76952 & 0.15955 \\
		\midrule
		Fisher & WIZ & 15 & 2014-09-23 & 2014-09-30 & 0.10010 & 0.01311 \\
		Fisher & WIZ & 68 & 2015-08-07 & 2015-08-14 & 0.09764 & 0.01285 \\
		Fisher & WIZ & 255 & 2018-09-02 & 2018-09-09 & 0.09643 & 0.01287 \\
		Fisher & WIZ & 330 & 2019-11-26 & 2019-12-03 & 0.10240 & 0.01298 \\
		\addlinespace[2pt]
		Fisher & WSRZ & 116 & 2016-05-21 & 2016-05-28 & 0.11279 & 0.01403 \\
		Fisher & WSRZ & 126 & 2016-07-20 & 2016-07-27 & 0.11506 & 0.01568 \\
		Fisher & WSRZ & 147 & 2016-11-23 & 2016-11-30 & 0.12750 & 0.02409 \\
		Fisher & WSRZ & 151 & 2016-12-17 & 2016-12-24 & 0.11120 & 0.02306 \\
		\bottomrule
	\end{tabular}
\end{table}

	
\subsection{Feature Correlations}
The seven-feature subset retains substantial Spearman correlation
among several descriptors Figures~\ref{fig:corr_wiz} and~\ref{fig:corr_wsrz} show that the seven retained features remain partly redundant at both stations. Strong positive correlations occur between Fisher information and Fisher-based complexity, and between Fisher-based complexity and disequilibrium. The variance of Shannon entropy is also strongly correlated with the variances of Shannon complexity and Tsallis entropy. These relationships indicate overlapping information among the retained descriptors and motivate further evaluation of reduced feature subsets.
	
\begin{figure}[H]
	\centering
	\includegraphics[width=\linewidth,height=0.71\textheight,keepaspectratio]
	{\figpath/Feature_Correlation_Spearman_WIZ.pdf}
	\caption{Spearman rank correlations among the seven WIZ features.}
	\label{fig:corr_wiz}
\end{figure}
	
\begin{figure}[H]
	\centering
	\includegraphics[width=\linewidth,height=0.71\textheight,keepaspectratio]
	{\figpath/Feature_Correlation_Spearman_WSRZ.pdf}
	\caption{Spearman rank correlations among the seven WSRZ features.}
	\label{fig:corr_wsrz}
\end{figure}
	
\subsection{DBSCAN Anomaly Results}
DBSCAN yields one non-noise cluster at each station: 293 clustered windows at WIZ and 287 at WSRZ. The remaining 42 (12.54\%) and 46 (13.81\%) are noise points. All anomalies belong to Classes 1--2; none of the three Class 3 windows is flagged at either station (Table~\ref{tab:dbscan}). Thus, this configuration identifies unusual feature vectors but misses every eruption-labeled window.
	
\begin{table}[H]
	\centering\small
	\caption{Corrected DBSCAN status by VAL class.}\label{tab:dbscan}
	\begin{tabularx}{\linewidth}{Ylrr}
		\toprule
 	Station & VAL class & Anomaly & Clustered\\
		\midrule
		WIZ & 1 & 41 & 279\\
		WIZ & 2 & 1 & 11\\
		WIZ & 3 & 0 & 3\\
		WSRZ & 1 & 43 & 275\\
		WSRZ & 2 & 3 & 9\\
		WSRZ & 3 & 0 & 3\\
		\bottomrule
	\end{tabularx}
\end{table}
	
\begin{table}[H]
	\centering\small
	\caption{DBSCAN binary performance: Class 3 positive; full-data escriptive evaluation.}\label{tab:db_metrics}
	\begin{tabularx}{\linewidth}{Yrrrrr}
		\toprule
		Station & Accuracy & Kappa & BA & Sensitivity & Specificity\\
		\midrule
		WIZ & 0.8657 & -0.0170 & 0.4367 & 0.0000 & 0.8735\\
		WSRZ & 0.8529 & -0.0172 & 0.4303 & 0.0000 & 0.8606\\
		\bottomrule
	\end{tabularx}
\end{table}
	
Binary eruption precision and F1 are zero for both stations. The non-eruption majority drives high accuracy and does not indicate successful eruption recognition. The 42 and 46 anomalies outside Class 3 are candidates for follow-up, not evidence of previously unrecorded eruptions.
	
The PCA displays preserve 88.1\% of total standardized-feature variance for WIZ (PC1 64.0\%, PC2 24.1\%) and 90.5\% for WSRZ (67.3\%, 23.2\%). DBSCAN itself operates in seven dimensions; PCA is only a display. Sankey flows count weekly windows and reproduce the cross-tabulation.
	
\begin{figure}[H]
	\centering
	\includegraphics[width=\linewidth,height=0.71\textheight,keepaspectratio]{\figpath/DBSCAN_VAL_PCA_plot_WIZ.pdf}
	\caption{WIZ: PCA display of the seven-feature data, coloured by VAL class and shaped by DBSCAN state.}\label{fig:dbscan_WIZ}
\end{figure}

\begin{figure}[H]
	\centering
	\includegraphics[width=\linewidth,height=0.71\textheight,keepaspectratio]{\figpath/DBSCAN_VAL_PCA_plot_WSRZ.pdf}
	\caption{WSRZ: PCA display of the seven-feature data, coloured by VAL class and shaped by DBSCAN state.}\label{fig:dbscan_WSRZ}
\end{figure}	
	
\begin{figure}[H]
	\centering
	\includegraphics[width=\linewidth,height=0.71\textheight,keepaspectratio]{\figpath/VAL_DBSCAN_Sankey_WIZ.pdf}
	\caption{WIZ: correspondence between VAL activity class and DBSCAN anomaly/clustered status. Flow widths count weekly windows.}\label{fig:sankey_WIZ}
\end{figure}

\begin{figure}[H]
	\centering
	\includegraphics[width=\linewidth,height=0.71\textheight,keepaspectratio]{\figpath/VAL_DBSCAN_Sankey_WSRZ.pdf}
	\caption{WSRZ: correspondence between VAL activity class and DBSCAN anomaly/clustered status. Flow widths count weekly windows.}\label{fig:sankey_WSRZ}
\end{figure}	
	
\subsection{Weighted SVM Results}
The selected parameters are cost $=10$, $\gamma=0.05$ for WIZ and cost $=100$, $\gamma=0.001$ for WSRZ. Table~\ref{tab:svm_cv} reports the repeated-validation summaries used for selection. WSRZ has higher mean one-versus-rest balanced accuracy, but lower overall accuracy and a similar macro F1. This does not establish uniformly superior classification at WSRZ.

\begin{table}[H]
	\centering\small
	\caption{SVM three-class tuning summaries: means over ten repeats of three-fold validation. BA is mean one-versus-rest balanced accuracy.}\label{tab:svm_cv}
	\begin{tabularx}{\linewidth}{Yrrrrr}
		\toprule
		Station & Accuracy & Kappa & BA & SD(BA) & Macro F1\\
		\midrule
		WIZ & 0.6042 & 0.0332 & 0.5380 & 0.0380 & 0.2873\\
		WSRZ & 0.5429 & 0.0300 & 0.6421 & 0.0709 & 0.2858\\
		\bottomrule
	\end{tabularx}
\end{table}

Mean macro sensitivities are 0.3686 (WIZ) and 0.5596 (WSRZ); mean macro specificities are 0.7075 and 0.7246. These are different from the single-pass metrics below. At WSRZ, the selected $\gamma$ lies at the lower grid boundary. Its mean BA of 0.642122 is very close to 0.641945 for cost $=100$, $\gamma=0.005$, so we cannot establish a uniquely optimal parameter choice.

Table~\ref{tab:svm_confusion} gives the separate single-pass out-of-fold matrices. Their accuracies are $191/335=0.5701$ and $165/333=0.4955$, and kappas are 0.0253 and 0.0133, respectively. They are not intended to reproduce the repeated means in Table~\ref{tab:svm_cv}.

\begin{table}[H]
	\centering\small
	\caption{Single three-fold out-of-fold SVM confusion matrices. Rows: prediction; columns: VAL reference class.}\label{tab:svm_confusion}
	\begin{tabularx}{\linewidth}{Ylrrr}
		\toprule
		Station & Predicted class & Ref. 1 & Ref. 2 & Ref. 3\\
		\midrule
		WIZ & 1 & 185 & 5 & 2\\
		WIZ & 2 & 89 & 6 & 1\\
		WIZ & 3 & 46 & 1 & 0\\
		WSRZ & 1 & 159 & 5 & 1\\
		WSRZ & 2 & 129 & 5 & 1\\
		WSRZ & 3 & 30 & 2 & 1\\
		\bottomrule
	\end{tabularx}
\end{table}

In this pass, WIZ misses all three eruption windows and predicts Class 3 for 47 non-eruption windows. WSRZ detects one eruption window, misses two, and predicts Class 3 for 32 non-eruption windows. Its eruption recall is $1/3=0.3333$, precision is $1/33=0.0303$, and F1 is 0.0556. The station-specific SVM projections show predictions by color and reference classes by shape; they do not show an independent test set.

\begin{figure}[H]
	\centering
	\includegraphics[width=\linewidth,height=0.71\textheight,keepaspectratio]{\figpath/SVM_OOF_Predictions_PCA_WIZ.pdf}
	\caption{WIZ: single-pass out-of-fold SVM predictions displayed in PCA coordinates. Colors denote predictions and shapes denote VAL-derived reference classes.}\label{fig:svm_WIZ}
\end{figure}

\begin{figure}[H]
	\centering
	\includegraphics[width=\linewidth,height=0.71\textheight,keepaspectratio]{\figpath/SVM_OOF_Predictions_PCA_WSRZ.pdf}
	\caption{WSRZ: single-pass out-of-fold SVM predictions displayed in PCA coordinates. Colors denote predictions and shapes denote VAL-derived reference classes.}\label{fig:svm_WSRZ}
\end{figure}
	

	
\section{Discussion}
The corrected results show that density anomalies and VAL eruption classes describe different aspects of these data. DBSCAN finds unusual windows but no Class 3 windows. These anomalies may reflect volcanic or environmental variation or data-quality effects, and require independent examination before assigning physical causes. The entropy-complexity figures provide useful descriptive context, but their labeled extremes cannot be interpreted as precursors without date matching and event timing.
	
SVM performance remains limited. The WSRZ repeated mean balanced accuracy exceeds WIZ's, and the separate WSRZ pass recovers one eruption window, but precision is low and kappa is close to zero. Only three eruption-labeled windows at each station contribute to these comparisons. Their count is not necessarily the number of independent eruption episodes. Class weighting cannot supply additional independent examples.
	
The seven-feature subset retains substantial Spearman correlation
among several descriptors. Station comparisons should also account for missing windows and differences in signal properties rather than attributing every difference to physical sensitivity.
	
The reported SVM scores are internal tuning and out-of-fold summaries, not external validation. A future evaluation should separate tuning from assessment and, if source windows overlap, prevent shared observations across folds. Recomputing class weights within each training partition would also avoid relying on full-data class frequencies. The clustering script does not specify the upstream window step or define the four variance features; those details remain to be reconciled with the feature-generation code.
	
A forecasting interpretation would require windows ending before an event, labels appropriate to a future prediction horizon, and evaluation based only on information available at the time. The present results support retrospective feature analysis and candidate-anomaly inspection, but not operational forecasting.
	
\section{Conclusion}
The corrected analysis uses 335 WIZ and 333 WSRZ weekly windows, with three eruption-labeled windows per station. DBSCAN identifies 42 and 46 anomalies, respectively, but none are eruption-labeled. Weighted SVM yields mean one-versus-rest balanced accuracies of 0.538 and 0.642 under repeated three-fold validation. In a separate out-of-fold pass, WIZ detects zero of three eruption windows, and WSRZ detects one of three, with many false-positive eruption predictions. The supplied plots and corrected tables therefore support exploratory characterization of RSAM feature variability, while reliable eruption recognition remains unestablished.
	
	\bibliographystyle{plain}
	\bibliography{seismic_references}
\end{document}